\documentclass[12pt]{article}
\usepackage[T1]{fontenc}
\usepackage[utf8]{inputenc}
\usepackage{lmodern}
\usepackage{textcomp}
\usepackage{enumitem}
\usepackage{geometry}
\geometry{margin=1in}
\begin{document}
\begin{center}
Answer Key - Religious Studies - K-12 Level Assessment 15 \
\small (Answers are ordered exactly as in the question paper.)
\end{center}

\section*{Multiple Choice}
\begin{enumerate}[label=\arabic*.]
\item \textbf{Answer:} C
\\ \textit{Options:} A) Early Third Century BCE; B) Second and First Century BCE; C) Late Sixth Century BCE; D) Fourth and Third Century BCE
\\ \textit{Note:} The late sixth century BCE marks a significant philosophical transition in ancient Greece when thinkers like Thales and Anaximander began to propose naturalistic explanations for the cosmos, moving away from traditional anthropomorphic representations of the divine.
\item \textbf{Answer:} D
\\ \textit{Options:} A) Soul; B) Sound; C) Non-matter; D) Non-soul
\\ \textit{Note:} In Jaina traditions, the term ajiva refers to all that is non-soul or non-living, distinguishing it from jiva, which signifies the soul or living entities.
\item \textbf{Answer:} A
\\ \textit{Options:} A) Enlightened being; B) Worthy ones; C) Saintly ones; D) Sages
\\ \textit{Note:} The term "Arhats" refers to enlightened beings in Buddhism who have attained Nirvana and freed themselves from the cycle of rebirth and suffering.
\item \textbf{Answer:} B
\\ \textit{Options:} A) Exactly 17; B) At least 17; C) At least 13; D) Exactly 24
\\ \textit{Note:} The Fatihah, which is the opening chapter of the Quran, is recited in every unit of the five daily prayers performed by Muslims, totaling a minimum of 17 recitations each day.
\item \textbf{Answer:} B
\\ \textit{Options:} A) Epistemology; B) Morality; C) Religion; D) Aesthetics
\\ \textit{Note:} In his final work, Laws, Plato emphasized the importance of morality as a guiding principle for human behavior and societal organization, moving away from the abstract principles of cosmology that characterized his earlier dialogues.
\end{enumerate}

\section*{True / False}
\begin{enumerate}[label=\arabic*.]
\setcounter{enumi}{5}
\item \textbf{Answer:} False
\\ \textit{Options:} A) True; B) False
\\ \textit{Note:} Guru Nanak began preaching the message of the divine Name at the age of 30, not 40.
\item \textbf{Answer:} True
\\ \textit{Options:} A) True; B) False
\\ \textit{Note:} Augustine was a prominent theologian who opposed Pelagius by asserting that original sin affected human nature and that divine grace is necessary for salvation, thus emphasizing the limitations of free will without God's intervention.
\item \textbf{Answer:} False
\\ \textit{Options:} A) True; B) False
\\ \textit{Note:} The statement is false because Moses is not historically recorded as having visited the Ka'ba, which is primarily associated with Islamic tradition and Abraham, whereas Moses is a central figure in Judaism and Christianity.
\item \textbf{Answer:} True
\\ \textit{Options:} A) True; B) False
\\ \textit{Note:} The statement is true because the Conservative branch of Judaism, established by Zacharias Frankel in the 19 th century, emphasizes a balanced approach to Jewish law and tradition, which is characterized by its understanding of Judaism as a living faith that evolves while maintaining historical continuity.
\item \textbf{Answer:} False
\\ \textit{Options:} A) True; B) False
\\ \textit{Note:} The Isma'ilis are primarily associated with the Shia branch of Islam, as they follow a distinct interpretation of Islamic teachings and leadership that diverges from Sunni beliefs.
\end{enumerate}

\section*{Long Form Response}
\begin{enumerate}[label=\arabic*.]
\setcounter{enumi}{10}
\item \textbf{Answer:} Mappo is a concept in Japanese Buddhism that signifies an age of decline, where the teachings of the Buddha are believed to lose their effectiveness, and moral decay spreads among people. This idea reflects broader themes in religious beliefs, such as the cyclical nature of time and the tension between spiritual ideals and human behavior. In this age, practitioners may feel a sense of hopelessness or despair, as they struggle to maintain their faith in a world perceived as increasingly corrupt. Mappo encourages individuals to seek personal enlightenment and deeper understanding, even in challenging times, highlighting the resilience of faith amidst decline. Ultimately, it serves as a reminder of the importance of spiritual practice and community support in navigating life's difficulties.
\\ \textit{Note:} The answer correctly identifies mappo as an age of decline in Japanese Buddhism, emphasizing its implications for practitioners' spiritual struggles and its reflection of broader religious themes such as moral decay and the cyclical nature of time.
\item \textbf{Answer:} The Daoist concept of wuwei, which translates to "non-action" or "effortless action," is a fundamental principle in Daoism that emphasizes aligning oneself with the natural flow of the universe. Instead of forcing actions or striving for control, wuwei encourages individuals to act spontaneously and harmoniously with their surroundings. This idea influences one's approach to life by promoting a sense of peace and balance, allowing people to respond to situations without unnecessary stress or effort. In nature, wuwei teaches respect for the natural world and its rhythms, suggesting that by observing and adapting to nature, rather than trying to dominate it, one can achieve greater harmony and fulfillment. Overall, wuwei highlights the importance of simplicity, patience, and a deep connection to the Dao, or the way of the universe.
\\ \textit{Note:} correct because it accurately describes wuwei as a key Daoist principle that advocates for effortless action in harmony with nature, highlighting its significance in promoting peace and balance in one's life.
\end{enumerate}

\vfill
\noindent\textit{End of Answer Key}
\end{document}